\chapter{Introduction}

Snow plays a critical role in Earth energy balance, 

\section{Thesis outline}

This thesis is structured as follows:
\begin{itemize}[label=$\blacktriangleright$]
    \item \textit{Chapter \ref{cap:snow_basics}} provides the theoretical background needed to understand the work presented in this thesis. The 
    first part focuses on snow physics. It begins with the nucleation of ice crystals in cold clouds and follows the snowflakes down to the 
    ground, where they accumulate to form a new layer. It then describes snow metamorphism and the energy balance of the snowpack, and finally 
    liquid water percolation and runoff generation. The second part reviews snow measurement techniques, distinguishing between manual observations, 
    automatic monitoring networks and remote sensing, and discussing the strengths and limitations of each approach. The third part examines how 
    climate change is affecting snow cover, outlining the current state of snow dynamics worldwide.
    \item \textit{Chapter \ref{cap:snow_products}} introduces the study area and the three main snow products that cover it. It opens with a 
    description of the Italian hydrological portion of the Greater Alpine Region (GAR), which encompasses the Italian Alps and the Northern Apennines. 
    Each dataset is examined in turn, with particular attention to how it is generated and the estimates it provides over Italy. The chapter closes 
    with a comparison of the three snow products, highlighting the discrepancies between their estimates and thereby motivating the need for a new 
    approach to snow monitoring at the national scale.
    \item \textit{Chapter \ref{cap:swe_model}} describes the model used to generate the snow water equivalent (SWE) dataset. The chapter is 
    structured to closely follow the modelling workflow. First, the precipitation and air temperature datasets used as model inputs are introduced, 
    focusing on the numerical methods employed to generate them and discussing some preliminary results, as they may significantly influence the 
    spatial and temporal distribution of snow cover. The precipitation phase partitioning scheme and the approach used to estimate snowmelt are then 
    described. The chapter closes with the description of the procedure used to generate the SWE maps.
    \item \textit{Chapter \ref{cap:swe_calibration}}
\end{itemize}